\chapter{Introduction}

\section{Motivation}
Terrestrial networks cover only a fraction of the Earth's land area, and cloud
data centres are usually far from user devices. Mega-constellations of LEO
satellites, combined with high-speed inter-satellite links (ISL), make it
possible to deploy \emph{Satellite Edge Nodes} (SENs) that offer low-latency
computing almost everywhere \citep{icdcs_sc,magliarisi2025orbit}.

SENs are constrained: CPU, memory and storage are scarce, the energy budget can
only be recharged in sunlight, and the high orbital speed means a satellite
remains visible to a ground user for a limited time (25--176\,s for an elevation
threshold of $40^\circ$ \citep{magliarisi2025orbit}). A high load therefore
leads quickly to congestion, longer delays and a lower success rate. Task
scheduling is crucial to mitigate these limits.

\section{Problem and Proposed Approach}
Existing heuristics (FCFS, HEFT, HADRT, \dots) work in specific scenarios but
take significant time and domain knowledge to design. Deep reinforcement
learning is an alternative but is costly to train and hard to interpret. This
thesis follows the \emph{genetic programming hyper-heuristic} (GPHH) idea: let
evolution discover a scheduling rule automatically, as a human-readable
expression over task and satellite features.

\section{Contributions}
\begin{itemize}
  \item A formulation of satellite task scheduling as a GP problem, with an
        expression-tree representation of the scheduling rule.
  \item A greedy decoder that turns a tree into a feasible schedule under
        orbital-window and energy constraints.
  \item A multi-objective fitness combining completion ratio, response time and energy.
  \item A parallel C++ implementation with Python bindings, together with a
        performance analysis (Chapter~\ref{ch:impl}).
  \item An experimental evaluation against hand-written heuristics, including a parameter study (Chapter~\ref{ch:exp}).
\end{itemize}

\section{Thesis Outline}
Chapter~\ref{ch:background} introduces genetic programming.
Chapter~\ref{ch:related} reviews GP-based scheduling and the SEC scheduling
literature. Chapter~\ref{ch:problem} formalises the problem.
Chapter~\ref{ch:method} describes the algorithm and Chapter~\ref{ch:impl} its
implementation. Chapter~\ref{ch:exp} reports the experiments and
Chapter~\ref{ch:concl} concludes.
